\documentclass[a4paper,12pt]{article}

\usepackage[T2A]{fontenc}
\usepackage[utf8]{inputenc}
\usepackage[russian]{babel}
\usepackage{amsmath,amssymb,amsthm,mathtools}
\usepackage{helvet}
\renewcommand{\familydefault}{\sfdefault}
\usepackage{icomma}
\usepackage{geometry}
\usepackage{microtype}
\usepackage{tikz}
\usepackage{tkz-euclide}
\usepackage{pgfplots}
\pgfplotsset{compat=1.18}
\usepackage{tabularx,booktabs,longtable}
\usepackage{enumitem}
\usepackage{hyperref}
\usepackage{parskip}
\usepackage{fancyhdr}
\usepackage{setspace}
\usepackage{xcolor}

\geometry{margin=2.1cm}
\onehalfspacing

\definecolor{accent}{HTML}{008080}
\definecolor{darkaccent}{HTML}{004D4D}
\definecolor{textgray}{HTML}{333333}
\definecolor{softgray}{HTML}{E9EEEE}
\definecolor{muted}{HTML}{667777}

\hypersetup{
    colorlinks=true,
    linkcolor=darkaccent,
    urlcolor=darkaccent
}

\setlength{\parindent}{0pt}
\setlength{\parskip}{0.55em}

\pagestyle{fancy}
\fancyhf{}
\lhead{\small\color{textgray} Теория вероятностей}
\rhead{\small\color{textgray} Метод позиций}
\cfoot{\small\color{muted}\thepage}
\renewcommand{\headrulewidth}{0.4pt}

\newcommand{\sectionline}{%
    \vspace{0.2em}
    {\color{accent}\rule{\linewidth}{0.7pt}}
    \vspace{0.4em}
}

\newcommand{\positionbox}[1]{%
    \tikz[baseline=-0.7ex]\node[
        draw=accent,
        rounded corners=2pt,
        minimum width=1.25cm,
        minimum height=0.8cm,
        line width=0.8pt
    ]{\textbf{#1}};%
}

\newcommand{\important}[1]{%
    \par\smallskip
    {\color{darkaccent}\textbf{#1}}
    \par\smallskip
}

\begin{document}

% =========================================================
% ТИТУЛЬНЫЙ ЛИСТ
% =========================================================

\begin{titlepage}
    \thispagestyle{empty}
    \centering

    \vspace*{2.8cm}

    {\Large\color{muted} ЧЕК-ЛИСТ\par}

    \vspace{0.8cm}

    {\fontsize{26}{32}\selectfont
    \bfseries\color{darkaccent}
    Теория вероятностей\par}

    \vspace{0.35cm}

    {\LARGE\bfseries
    Метод позиций и подсчёт исходов\par}

    \vspace{1.4cm}

    {\large
    Как быстро считать количество комбинаций\\
    без громоздкого полного перебора\par}

    \vfill

    {\large
    Лёвин Артём Александрович\\[0.25em]
    эксклюзивно для Ксении Клыковой\par}

    \vspace{0.8cm}

    {\large \today\par}

    \vspace{1.4cm}

    \begin{tikzpicture}[x=1cm,y=1cm]
        \draw[
            color=accent!60,
            line width=1.2pt
        ]
        (0,0)
        .. controls (3,0.75) and (7,-0.75) ..
        (10,0)
        .. controls (12,0.45) and (14,0.35) ..
        (15,0.1);
    \end{tikzpicture}

    \vspace{0.8cm}
\end{titlepage}

% =========================================================
% ОСНОВНОЙ ЧЕК-ЛИСТ
% =========================================================

\section*{1. Главная формула классической вероятности}
\sectionline

Если все элементарные исходы равновозможны, то

\[
\boxed{
P(A)=\frac{N(A)}{N}
}
\]

где

\begin{itemize}[leftmargin=1.5em]
    \item \(N\) --- общее количество возможных исходов;
    \item \(N(A)\) --- количество исходов, благоприятных событию \(A\).
\end{itemize}

\important{Главная задача при решении:}
сначала точно понять условие события, затем корректно посчитать общее и благоприятное количество исходов.

\subsection*{Перевод текста на математический язык}

В задачах по вероятности формулировка часто важнее самих вычислений.

\begin{center}
\renewcommand{\arraystretch}{1.3}
\begin{tabularx}{0.92\linewidth}{@{}lX@{}}
\toprule
Формулировка & Математический смысл \\
\midrule
меньше \(4\) & \(x<4\) \\
не более \(4\) & \(x\le 4\) \\
не менее \(4\) & \(x\ge 4\) \\
больше \(4\) & \(x>4\) \\
отличается от \(3\) на единицу & \(|x-3|=1\) \\
ровно три раза & событие произошло в точности в трёх испытаниях \\
\bottomrule
\end{tabularx}
\end{center}

\subsection*{Пример 1. Один бросок кубика}

Кубик бросают один раз. Найдём вероятность выпадения числа меньше \(4\).

Подходящие результаты:

\[
1,\ 2,\ 3.
\]

Следовательно,

\[
P=\frac{3}{6}=\frac12.
\]

Если бы требовалось число \emph{не более \(4\)}, подходили бы

\[
1,\ 2,\ 3,\ 4,
\]

поэтому

\[
P=\frac46=\frac23.
\]

\subsection*{Пример 2. Отличие на единицу}

Требуется число, отличающееся от \(3\) на единицу:

\[
|x-3|=1.
\]

Получаем два подходящих значения:

\[
x=2,\qquad x=4.
\]

Поэтому

\[
P=\frac26=\frac13.
\]

\newpage

\section*{2. Метод позиций}
\sectionline

Метод позиций помогает считать количество исходов через последовательное заполнение нескольких мест.

Каждая позиция отвечает за определённую часть результата:

\begin{itemize}[leftmargin=1.5em]
    \item результат первого броска;
    \item результат второго броска;
    \item место человека в очереди;
    \item первый выбранный участник;
    \item второй выбранный участник;
    \item обладатель первого приза;
    \item обладатель второго приза;
    \item и т.\,д.
\end{itemize}

Число внутри позиции означает \textbf{количество допустимых вариантов для этой позиции}.

\[
\positionbox{6}
\qquad
\positionbox{6}
\]

означает:

\begin{itemize}[leftmargin=1.5em]
    \item для первой позиции имеется \(6\) вариантов;
    \item для второй позиции имеется \(6\) вариантов.
\end{itemize}

Количество комбинаций:

\[
6\cdot6=36.
\]

\important{Ключевой вопрос:}
сколько вариантов можно поставить на каждую позицию с учётом уже сделанных выборов и условия задачи?

\subsection*{Алгоритм метода позиций}

\begin{enumerate}[leftmargin=1.7em]
    \item Определите, что представляет собой одна позиция.
    \item Определите количество позиций.
    \item Для каждой позиции посчитайте количество допустимых вариантов.
    \item Перемножьте количества вариантов.
    \item Отдельно посчитайте общее и благоприятное количество исходов.
    \item Подставьте их в формулу вероятности.
\end{enumerate}

\[
\boxed{
N=k_1k_2\cdots k_n
}
\]

где \(k_i\) --- количество допустимых вариантов на \(i\)-й позиции.

\section*{3. Два броска кубика}
\sectionline

Кубик бросают дважды.

Для каждого броска возможно \(6\) результатов:

\[
\positionbox{6}
\qquad
\positionbox{6}.
\]

Следовательно,

\[
N=6\cdot6=36.
\]

\subsection*{Одинаковые числа}

Для первого броска подходит любое из \(6\) чисел.

После его результата второй бросок должен совпасть с первым:

\[
\positionbox{6}
\qquad
\positionbox{1}.
\]

Поэтому

\[
N(A)=6\cdot1=6,
\]

а вероятность равна

\[
\boxed{
P(A)=\frac6{36}=\frac16
}.
\]

\subsection*{Различные числа}

Первый результат произволен: \(6\) вариантов.

Для второго результата подходят все числа, кроме уже выпавшего:

\[
\positionbox{6}
\qquad
\positionbox{5}.
\]

Получаем

\[
N(A)=6\cdot5=30,
\]

\[
\boxed{
P(A)=\frac{30}{36}=\frac56
}.
\]

\newpage

\section*{4. Три броска кубика}
\sectionline

Общее количество исходов:

\[
\positionbox{6}
\qquad
\positionbox{6}
\qquad
\positionbox{6},
\]

\[
N=6^3=216.
\]

\subsection*{Все три числа одинаковы}

\[
\positionbox{6}
\qquad
\positionbox{1}
\qquad
\positionbox{1}.
\]

Следовательно,

\[
N(A)=6,
\]

\[
\boxed{
P(A)=\frac6{216}=\frac1{36}
}.
\]

\subsection*{Все три числа различны}

\[
\positionbox{6}
\qquad
\positionbox{5}
\qquad
\positionbox{4}.
\]

Следовательно,

\[
N(A)=6\cdot5\cdot4=120,
\]

\[
\boxed{
P(A)=\frac{120}{216}=\frac59
}.
\]

\subsection*{Ровно два числа одинаковы}

Все \(216\) исходов распадаются на три группы:

\begin{enumerate}[leftmargin=1.7em]
    \item все три числа одинаковы;
    \item все три числа различны;
    \item ровно два числа одинаковы.
\end{enumerate}

Поэтому

\[
N(\text{ровно два одинаковы})
=
216-6-120=90.
\]

Следовательно,

\[
\boxed{
P=\frac{90}{216}=\frac5{12}
}.
\]

Этот результат можно получить и прямым подсчётом:

\[
6\cdot5\cdot3=90.
\]

Здесь:

\begin{itemize}[leftmargin=1.5em]
    \item \(6\) --- выбор повторяющегося числа;
    \item \(5\) --- выбор отличающегося числа;
    \item \(3\) --- выбор позиции, на которой стоит отличающееся число.
\end{itemize}

\section*{5. Четыре броска кубика}
\sectionline

Общее количество комбинаций:

\[
N=6^4=1296.
\]

\subsection*{Все четыре числа одинаковы}

\[
N(A)=6\cdot1\cdot1\cdot1=6.
\]

Поэтому

\[
\boxed{
P(A)=\frac6{6^4}=\frac1{216}
}.
\]

\subsection*{Все четыре числа различны}

\[
N(A)=6\cdot5\cdot4\cdot3=360.
\]

Поэтому

\[
\boxed{
P(A)=\frac{360}{1296}=\frac5{18}
}.
\]

\newpage

\section*{6. «И» и «или»: важное уточнение}
\sectionline

В задачах удобно мысленно связывать этапы словами «и» и «или», однако использовать этот приём нужно строго.

\subsection*{Последовательные позиции}

Если результат строится последовательным выбором:

\[
\text{первая позиция И вторая позиция И третья позиция},
\]

то количества вариантов перемножаются:

\[
N=k_1k_2k_3.
\]

Например, три различных результата кубика:

\[
6\cdot5\cdot4.
\]

\subsection*{Альтернативные несовместимые случаи}

Если благоприятное событие состоит из нескольких \textbf{непересекающихся} вариантов, их количества складываются.

Например:

\[
\text{все числа одинаковы ИЛИ все числа различны}.
\]

Эти два случая одновременно произойти не могут, поэтому

\[
N=6+120=126.
\]

\important{Осторожно с вероятностями.}

Для событий \(A\) и \(B\)

\[
P(A\cup B)=P(A)+P(B)
\]

верно непосредственно только для несовместимых событий.

В общем случае:

\[
\boxed{
P(A\cup B)=P(A)+P(B)-P(A\cap B)
}.
\]

А формула

\[
P(A\cap B)=P(A)P(B)
\]

в таком виде применяется для независимых событий.

В общем случае используется

\[
\boxed{
P(A\cap B)=P(A)\,P(B\mid A)
}.
\]

\section*{7. Очередь}
\sectionline

Андрей, Борис и Виктор случайно становятся в очередь.

Найдём вероятность того, что Виктор окажется раньше Андрея.

\subsection*{Общее количество порядков}

\[
\positionbox{3}
\qquad
\positionbox{2}
\qquad
\positionbox{1}.
\]

Следовательно,

\[
N=3\cdot2\cdot1=6.
\]

Ровно в половине перестановок Виктор расположен раньше Андрея:

\[
N(A)=3.
\]

Поэтому

\[
\boxed{
P(A)=\frac36=\frac12
}.
\]

\important{Типичная ошибка:}
при ручном переборе легко включить комбинацию, в которой порядок двух нужных людей нарушает условие. После составления списка всегда проверяйте каждую комбинацию непосредственно по тексту задачи.

\section*{8. Выбор людей: ловушка порядка}
\sectionline

Пусть из группы из \(5\) туристов случайно выбирают двух человек. Требуется вероятность того, что турист \(A\) попадёт в выбранную пару.

Для применения метода позиций удобно временно различать первого и второго выбранного туриста.

Общее количество упорядоченных исходов:

\[
\positionbox{5}
\qquad
\positionbox{4},
\]

\[
N=5\cdot4=20.
\]

Турист \(A\) может оказаться:

\begin{itemize}[leftmargin=1.5em]
    \item на первой позиции: \(1\cdot4=4\) исхода;
    \item на второй позиции: \(4\cdot1=4\) исхода.
\end{itemize}

Следовательно,

\[
N(A)=4+4=8,
\]

\[
\boxed{
P(A)=\frac8{20}=\frac25
}.
\]

\important{Очень важная проверка.}
Если задача спрашивает, попадёт ли человек в пару, необходимо учитывать \textbf{все позиции}, на которых он может оказаться.

Для группы из \(n\) человек при выборе двух:

\[
P(A)=\frac{2}{n}.
\]

Например, для \(500\) человек:

\[
\boxed{
P(A)=\frac2{500}=\frac1{250}
}.
\]

\newpage

\section*{9. Геометрическая задача: диагональ квадрата}
\sectionline

Из четырёх вершин квадрата случайно выбирают две различные вершины. Найдём вероятность того, что выбранные вершины являются концами диагонали.

\begin{center}
\begin{tikzpicture}[scale=1.05]
    \tkzDefPoint(0,0){A}
    \tkzDefPoint(4,0){B}
    \tkzDefPoint(4,4){C}
    \tkzDefPoint(0,4){D}

    \tkzDrawPolygon[line width=1pt,color=accent](A,B,C,D)
    \tkzDrawSegments[line width=1pt,color=muted](A,C B,D)
    \tkzDrawPoints(A,B,C,D)
    \tkzLabelPoints[below](A,B)
    \tkzLabelPoints[above](C,D)
\end{tikzpicture}
\end{center}

Для первой вершины имеется \(4\) варианта, для второй --- \(3\):

\[
N=4\cdot3=12.
\]

Если первая вершина уже выбрана, ровно одна вершина лежит напротив неё и образует диагональ:

\[
N(A)=4\cdot1=4.
\]

Следовательно,

\[
\boxed{
P(A)=\frac4{12}=\frac13
}.
\]

\section*{10. Монета и повторные испытания}
\sectionline

При каждом броске монеты имеются два результата:

\[
\text{О --- орёл},\qquad
\text{Р --- решка}.
\]

Если монету бросают \(n\) раз, то

\[
\boxed{
N=2^n
}.
\]

\subsection*{Элементарный исход ОРО}

Монету бросают три раза.

Общее количество исходов:

\[
N=2^3=8.
\]

Конкретный исход ОРО единственный:

\[
N(A)=1.
\]

Следовательно,

\[
\boxed{
P(\text{ОРО})=\frac18
}.
\]

\subsection*{Ровно три орла из четырёх бросков}

Общее количество исходов:

\[
N=2^4=16.
\]

Подходящие последовательности:

\[
\text{ОООР},\quad
\text{ООРО},\quad
\text{ОРОО},\quad
\text{РООО}.
\]

Их \(4\), поэтому

\[
\boxed{
P=\frac4{16}=\frac14
}.
\]

\section*{11. Три матча}
\sectionline

Перед каждым матчем жеребьёвка определяет, начнёт ли команда игру с мячом.

Для каждого матча имеется два равновозможных исхода:

\[
\text{начинает},\qquad
\text{не начинает}.
\]

Для трёх матчей:

\[
N=2^3=8.
\]

Требуется начать игру ровно два раза.

Подходят три варианта расположения двух успешных матчей:

\[
(1,2),\qquad
(1,3),\qquad
(2,3).
\]

Поэтому

\[
N(A)=3,
\]

\[
\boxed{
P(A)=\frac38
}.
\]

Здесь метод позиций эффективно даёт общее количество исходов, а небольшое количество благоприятных случаев удобно найти перебором.

\newpage

\section*{12. Когда использовать таблицу, а когда позиции}
\sectionline

\begin{center}
\renewcommand{\arraystretch}{1.35}
\begin{tabularx}{\linewidth}{@{}p{0.22\linewidth}XX@{}}
\toprule
Метод & Сильные стороны & Когда особенно полезен \\
\midrule
Таблица
&
Очень наглядна; позволяет увидеть все исходы
&
Два небольших множества вариантов; первый этап знакомства с задачей
\\
\addlinespace
Метод позиций
&
Быстрый; легко масштабируется; позволяет обходиться без длинного перебора
&
Несколько бросков, очереди, распределения, выбор участников и призов
\\
\addlinespace
Гибридный подход
&
Общее число исходов считается позициями, небольшое число благоприятных --- перебором
&
События вида «ровно два раза», «ровно три раза» при небольшом количестве испытаний
\\
\bottomrule
\end{tabularx}
\end{center}

\section*{13. Самопроверка решения}
\sectionline

Перед записью ответа проверьте:

\begin{enumerate}[leftmargin=1.7em]
    \item Точно ли расшифровано условие?
    \item Все ли исходы в знаменателе равновозможны?
    \item Что означает каждая позиция?
    \item Числа в позициях --- это количества вариантов?
    \item Изменяется ли количество вариантов после предыдущего выбора?
    \item Важен ли порядок?
    \item Если порядок введён искусственно, учтены ли все позиции благоприятного объекта?
    \item Случаи, которые складываются, действительно несовместимы?
    \item Благоприятное количество не превосходит общего?
    \item Полученная вероятность находится в пределах
    \[
    0\le P\le1?
    \]
\end{enumerate}

\section*{14. Ключевые модели занятия}
\sectionline

\begin{center}
\renewcommand{\arraystretch}{1.35}
\begin{tabularx}{\linewidth}{@{}XcX@{}}
\toprule
Ситуация & Количество исходов & Идея \\
\midrule
\(n\) бросков кубика & \(6^n\) & на каждом броске \(6\) вариантов \\
\(n\) бросков монеты & \(2^n\) & на каждом броске \(2\) варианта \\
\(n\) различных людей в очереди & \(n(n-1)\cdots1\) & после каждого выбора остаётся на одного меньше \\
3 броска, все числа одинаковы & \(6\) & \(6\cdot1\cdot1\) \\
3 броска, все различны & \(120\) & \(6\cdot5\cdot4\) \\
3 броска, ровно два одинаковы & \(90\) & \(6\cdot5\cdot3\) \\
4 броска, все одинаковы & \(6\) & \(6\cdot1\cdot1\cdot1\) \\
4 броска, все различны & \(360\) & \(6\cdot5\cdot4\cdot3\) \\
\bottomrule
\end{tabularx}
\end{center}

\important{Главная идея занятия:}
метод позиций превращает длинный перебор в последовательность вопросов:

\[
\boxed{
\text{«Сколько вариантов подходит на эту позицию?»}
}
\]

После каждого сделанного выбора количество допустимых претендентов может изменяться.

% =========================================================
% ТРЕНИРОВОЧНЫЙ БЛОК
% =========================================================

\newpage
\section*{Тренировочный блок}
\sectionline

\textbf{Путь А.} Выполните 10 заданий. Старайтесь применять метод позиций там, где он сокращает решение.

\subsection*{Средний уровень}

\begin{enumerate}[leftmargin=2em,label=\textbf{\arabic*.}]

    \item
    Игральный кубик бросают один раз. Найдите вероятность того, что выпавшее число не менее \(4\).

    \item
    Игральный кубик бросают дважды. Найдите вероятность того, что оба раза выпадет одинаковое число очков.

    \item
    Монету бросают четыре раза. Найдите вероятность элементарного исхода
    \[
    \text{ОРОР}.
    \]

\end{enumerate}

\subsection*{Выше среднего}

\begin{enumerate}[leftmargin=2em,label=\textbf{\arabic*.},resume]

    \item
    Игральный кубик бросают трижды. Найдите вероятность того, что все три выпавших числа различны.

    \item
    Игральный кубик бросают трижды. Найдите вероятность того, что ровно два из трёх выпавших чисел одинаковы.

    \item
    Андрей, Борис, Виктор и Григорий случайно становятся в очередь. Найдите вероятность того, что Андрей окажется раньше Бориса.

    \item
    В группе \(8\) туристов случайно выбирают двух человек. Найдите вероятность того, что турист \(A\) попадёт в выбранную пару.

    \item
    Из четырёх вершин квадрата случайно выбирают две различные вершины. Найдите вероятность того, что они являются концами стороны квадрата.

    \item
    Монету бросают пять раз. Найдите вероятность того, что орёл выпадет ровно четыре раза.

\end{enumerate}

\subsection*{Сложный уровень}

\begin{enumerate}[leftmargin=2em,label=\textbf{\arabic*.},resume]

    \item
    Игральный кубик бросают четыре раза. Найдите вероятность того, что среди выпавших чисел имеется ровно три одинаковых числа, а четвёртое отличается от них.

\end{enumerate}

% =========================================================
% ОТВЕТЫ
% =========================================================

\newpage
\section*{Ответы к тренировочному блоку}
\sectionline

\renewcommand{\arraystretch}{1.45}

\begin{center}
\begin{tabularx}{0.82\linewidth}{@{}cX@{}}
\toprule
\textbf{№} & \textbf{Ответ} \\
\midrule

1 &
\[
\frac{3}{6}=\frac12
\]
\\

2 &
\[
\frac{6}{36}=\frac16
\]
\\

3 &
\[
\frac1{16}
\]
\\

4 &
\[
\frac{6\cdot5\cdot4}{6^3}
=
\frac59
\]
\\

5 &
\[
\frac{6\cdot5\cdot3}{6^3}
=
\frac5{12}
\]
\\

6 &
\[
\frac12
\]
\\

7 &
\[
\frac{2}{8}=\frac14
\]
\\

8 &
\[
\frac{8}{12}=\frac23
\]
\\

9 &
\[
\frac{5}{32}
\]
\\

10 &
\[
\frac{6\cdot5\cdot4}{6^4}
=
\frac{120}{1296}
=
\frac5{54}
\]
\\

\bottomrule
\end{tabularx}
\end{center}

\vfill

\begin{center}
{\color{darkaccent}\rule{0.65\linewidth}{0.7pt}}

\vspace{0.5em}

\textbf{Финальная памятка}

\[
\boxed{
P(A)=\frac{\text{благоприятные исходы}}
{\text{все равновозможные исходы}}
}
\]

\vspace{0.4em}

Метод позиций:
\[
\boxed{
N=k_1\cdot k_2\cdot\ldots\cdot k_n
}
\]

Перед вычислением всегда определяйте,
\textbf{что означает каждая позиция и сколько вариантов на неё претендует}.
\end{center}

\end{document}